% A.63 reserved by the root.  No priority claim is made.
\subsection{The coefficient-prime Tate class and a transcendence obstruction}

\begin{statement}{A.63}{Proved}
{A dyadic completed coordinate is recognized without a multiplier restriction}
Let $s$ be a noncuspidal Selmer section in the finite-ramification
range A.14, and let $A,B\in\widehat{\Q^\times}$ be its Kummer
coordinates for $t,1-t$.  If either coordinate belongs to
\[
                  \{\epsilon 2^\beta:
                         \epsilon\in\{1,-1\},\quad
                         \beta\in\widehat\Z\},
\tag{A.63.1}
\]
then $s$ is a rational point section.

Consequently, every labeled integral realization $E_{\lambda_0}$
in A.16 with
\[
            \lambda_0\in\{\pm2^m:m\in\Z\}
       \quad\hbox{or}\quad
            1-\lambda_0\in\{\pm2^m:m\in\Z\}
\tag{A.63.2}
\]
is recognized: $s=s_{\lambda_0}$.  The same conclusion holds for
the six fractional-linear transforms of these parameters under
the permutations of $0,1,\infty$.
There is no parity condition on $m$ and no determinant condition
on the integral intertwiner.  In particular this strengthens
A.54 for every parameter in its statement.
\end{statement}
\begin{proof}
We first prove the local assertion which supplies the additional
information.  Let $E/\Q$ be potentially multiplicative at a prime
$p$.  Put $j_0=j(E)$ and let $q_0\in\Q_p^\times$ be its Tate
parameter, characterized by
\[
              0<|q_0|_p<1,\qquad J(q_0)=j_0,
       \qquad J(q)=q^{-1}+744+196884q+\cdots.
\tag{A.63.3}
\]
This parameter lies in $\Q_p$ before taking a splitting extension:
the series $1/J(q)\in q\Z[[q]]$ has linear coefficient one, so
its inverse belongs to $z\Z[[z]]$ and converges at
$z=1/j_0\in p\Z_p$.  After a quadratic extension the curve is
the split Tate curve with this parameter.

For every rational number $r\ne0$, one has
\[
                 \kappa_p(q_0)\notin
                 \Q_p\,\kappa_p(r)
               \quad\hbox{in }H^1(\Q_p,\Q_p(1)).
\tag{A.63.4}
\]
Here and below the Kummer map includes the map from integral to
rational coefficients.  To prove (A.63.4), suppose that
$\kappa_p(q_0)=a\kappa_p(r)$ for $a\in\Q_p$ and put
$N=v_p(q_0)>0$, $m=v_p(r)\in\Z$.
The valuation map on Kummer cohomology gives $N=am$.
If $m=0$, this is already impossible.  Otherwise $a=N/m$ and
\[
                       \kappa_p(q_0^m/r^N)=0.
\tag{A.63.5}
\]
The kernel of the rational Kummer map on $\Q_p^\times$ is its
finite group of roots of unity.  Indeed, decompose that group as
$p^\Z$ times its roots of unity times the principal units; after
tensoring with $\Q_p$, valuation detects the first factor and
the $p$-adic logarithm detects the last factor.  For $p=2$ use
$\Q_2^\times=2^\Z\times\{\pm1\}\times(1+4\Z_2)$.
Thus (A.63.5) makes $q_0^m/r^N$ a root of unity, and makes
$q_0$ algebraic over $\Q$.

The proved $p$-adic Mahler--Manin theorem of
Barr\'e-Sirieix, Diaz, Gramain and Philibert states that a
nonzero algebraic $q\in\mathbf C_p$ with $|q|_p<1$ has
transcendental $J(q)$
\cite{LegBDGP1996,LegGramain1999,LegGraftieaux2005}.
The precise $p$-adic statement is opened in the coauthor's
account \cite[pp.~73--74]{LegGramain1999} and independently in
\cite[p.~180]{LegGraftieaux2005}.
It contradicts $J(q_0)=j_0\in\Q$, proving (A.63.4).
The same assertion holds after restriction to any finite
$F/\Q_p$: on rational cohomology the composite of restriction
and corestriction is multiplication by $[F:\Q_p]$, so
restriction is injective.  Both classes in (A.63.4) are
restrictions of the indicated $\Q_p$-classes.

We next give the exact map from a cusp packet to this local
calculation.  At the cusps $0,1,\infty$ use respectively
\[
                         w=t,\quad w=1-t,\quad w=1/t.
\tag{A.63.6}
\]
The Legendre $j$ function is unchanged by these substitutions.
Formal inversion of (A.63.3) in
$\Q((w))$ gives the universal Tate parameter
\[
                 q_T=\frac{w^2}{256}F(w),
                       \qquad F(w)\in1+w\Q[[w]].
\tag{A.63.7}
\]
For every positive integer $n$, the unit $F(w)$ has a unique
$n$th root with constant term one in $\Q[[w]]$; these roots
are compatible under the power maps.  Its Kummer class on the
formal puncture $\operatorname{Spec}\Q_p((w))$ is therefore
zero, and remains zero on any cusp section pulled back through
its decomposition group.  This explains the specialization
of the formal unit; no value of a convergent power series at a
completed Kummer coordinate is being assigned.

At cusp zero the nodal slopes are $\pm\sqrt{-1}$, giving
the constant quadratic splitting character of $-1$.
At cusp one the nodal slopes are $\pm1$, giving the trivial
character.  At infinity the actual change of variables
\[
                       x=tX,\qquad y=t\sqrt t\,Y
\tag{A.63.8}
\]
identifies the family with the Legendre family of parameter
$1/t$, after the indicated quadratic extension; its additional
twist is the character of $t$.  Thus after a finite local
extension killing the indicated characters, the three pulled
back Tate extension classes are exactly
\[
\begin{array}{c|c|c}
 \text{cusp}&\text{completed Tate class}&
                \text{coefficient-}p\text{ extension class}\\ \hline
 0&A^2/2^8&2\kappa_p(A)-8\kappa_p(2)\\
 1&B^2/2^8&2\kappa_p(B)-8\kappa_p(2)\\
 \infty&A^{-2}/2^8&-2\kappa_p(A)-8\kappa_p(2).
\end{array}
\tag{A.63.9}
\]
In the last row $B/A=-1$ on the cusp packet, so it may
equivalently be written with $B^{-2}$.
These are the full extension classes, including the constant
$256$.

Let $s_p$ be a local cusp restriction of $s$.  A.14 gives
one elliptic curve $E/\Q$ realizing all its rational Tate
representations; A.37 proves that $E$ is potentially
multiplicative at $p$.  Choose a common finite extension
$F/\Q_p$ splitting the realized curve and the cusp character.
The actual $p$-adic Tate representation has the exact sequence
\[
       0\longrightarrow\Q_p(1)\longrightarrow V_p(E)
                       \longrightarrow\Q_p\longrightarrow0,
\tag{A.63.10}
\]
of class $\kappa_p(q_0)$.
Its cyclotomic line is unique: the composite of a homomorphism
from $\Q_p(1)$ with the trivial quotient is zero, since the
cyclotomic character on $G_F$ is not trivial.  It follows that
the realizing isomorphism, in these Tate bases, is triangular
with two nonzero diagonal entries $a,d\in\Q_p^\times$.
Comparing upper-right cocycles, and then their cohomology
classes, identifies the cusp extension class with
$(a/d)\kappa_p(q_0)$.
This argument works at the coefficient prime itself and
does not replace either representation by its semisimplification.
In particular a cusp extension class lying on
$\Q_p\kappa_p(r)$ for any $r\in\Q^\times$ contradicts
(A.63.4), including after the restriction to $F$ just used.

Suppose now that $A=\epsilon2^\beta$ as in (A.63.1).
The cusp-zero and cusp-infinity rows of (A.63.9) have classes
\[
        (2\beta_p-8)\kappa_p(2),\qquad
        (-2\beta_p-8)\kappa_p(2),
\tag{A.63.11}
\]
respectively.  The sign $\epsilon$ is killed by the square.
Both possibilities contradict (A.63.4).  Thus every local
cuspidal restriction has label one, and the one-label
recognition theorem A.37 makes $s$ a rational point section.
If $B=\epsilon2^\beta$, the same argument excludes label one
and infinity; all remaining cuspidal restrictions have label
zero and A.37 applies again.

For the application, take $\lambda_0=\epsilon2^m$ in (A.63.2).
At each odd prime $p$, this parameter is a unit.  Where
$1-\lambda_0$ is also a unit, the realized curve has good
reduction and the local comparison gives an ordinary integral
point.  At an odd prime dividing $1-\lambda_0$, the canonical
saturated cyclotomic line in $T_2$, reduced in the labeled
$E[2]$, is $\langle(0,0)\rangle$.  The exact line comparison
in A.41 then makes the local section a cusp-one section or
an ordinary point with $v_p(1-x_p)>0$.  In both cases
$v_p(A)=0$.  This argument depends only on
$v_p(\lambda_0)=0$ and is valid also when $m<0$ or
$\epsilon=-1$.  Thus $A$ has zero valuation at every odd
prime and is precisely an element of (A.63.1), by the sign
and valuation description of $\widehat{\Q^\times}$.
The preceding argument recognizes $s$.  If
$1-\lambda_0=\epsilon2^m$, interchange $A$ and $B$.
The labeled integral isomorphism then identifies the rational
point parameter with $\lambda_0$ by the primitive-isogeny
and ordered-two-torsion argument at the end of A.54; that
argument uses a general integral intertwiner and imposes
no determinant-one condition.

For completeness, the permutation claim is obtained by
actual maps, not by an identification of unrelated
representations.  The automorphisms $f(t)=1-t$ and
$h(t)=1/t$ generate the six permutations.
For $f$, the substitution
$(x,y)\mapsto(1-x,iy)$ identifies the two elliptic curves
over $\Q(i)$ and twists the Tate representation by the
constant character of $-1$.
For $h$, (A.63.8) twists it by the character of $t$.
On the section this last character is $\kappa_2(A)$, whose
mod-two class is $\kappa_2(\lambda_0)$ by A.16.
Thus the two twists on the section and its realized rational
point are the same, and the labeled integral realization
passes to $E_{f(\lambda_0)}$ or $E_{h(\lambda_0)}$.
The substitutions permute the displayed two-torsion points
in precisely the same way on both sides.
Apply the result just proved after this automorphism and
then its inverse to the resulting point section.
\end{proof}

\paragraph{The exact simultaneous equations with more than one support prime.}
The preceding obstruction also has a direct form for the
remaining finite-support problem.  Write
\[
         A=\epsilon_A\prod_{r\in S}r^{a_r},\qquad
         B=\epsilon_B\prod_{r\in S}r^{b_r},
            \qquad a_r,b_r\in\widehat\Z,
\]
and include $2$ in $S$, with zero exponent when necessary.
For a cusp label $c$ define completed exponent vectors
\[
 d_{0,r}=2a_r-8\mathbf1_{r=2},\quad
 d_{1,r}=2b_r-8\mathbf1_{r=2},\quad
 d_{\infty,r}=-2a_r-8\mathbf1_{r=2}.
\tag{A.63.12}
\]
At a local cusp of label $c$ and residue prime $p$, let
$q_{0,p}$ be the realized Tate parameter and
$N_p=v_p(q_{0,p})>0$.
Write $d_{c,r,p}$ for the $p$-component of $d_{c,r}$.
The actual coefficient-$p$ comparison gives, with one
nonzero scalar $\rho_p$,
\[
 \sum_{r\in S}d_{c,r,p}\kappa_p(r)
       =\rho_p\kappa_p(q_{0,p}),\qquad
 d_{c,p,p}=\rho_pN_p.
\tag{A.63.13}
\]
The convention is $d_{c,p,p}=0$ for $p\notin S$.
Taking logarithms and eliminating $\rho_p$ gives the exact
equation
\[
 N_p\sum_{r\in S\setminus\{p\}}
                d_{c,r,p}\log_p(r)
       =d_{c,p,p}\log_p(q_{0,p}).
\tag{A.63.14}
\]
For an integral intertwiner of Weil multiplier
$\delta=(\delta_\ell)_\ell$, comparison of the same Tate
bases at every coefficient prime gives
\[
       \frac{d_{c,p,\ell}}{N_p}
                    =\frac{u_{p,\ell}^{\,2}}{\delta_\ell},
                 \qquad u_{p,\ell}\in\Z_\ell^\times.
\tag{A.63.15}
\]
This retains the same global multiplier at all local primes.
It also proves that the ratios on the left at any two
cuspidal primes have quotient in
$(\Z_\ell^\times)^2$.  The ordinary-point case instead
has positive integral Tate valuations on both sides and
recovers A.34.

The regular-coordinate constraints remain exact Kummer
equations: $B_p=1$ at label zero, $A_p=1$ at label one,
and $(B/A)_p=-1$ at infinity.  They give their respective
valuation and logarithm equations in addition to
(A.63.13)--(A.63.15), and retain the finite unit components
before passing to rational coefficients.

The map in (A.63.13) explains precisely why the one-line
argument alone does not recognize an arbitrary larger
support.  For $p\ne2$ choose $u=2$, and for $p=2$ choose
$u=3$.  Then $\log_p(u)\ne0$, and
\[
 \Q_p^2\longrightarrow H^1(\Q_p,\Q_p(1)),\qquad
       (x,y)\longmapsto x\kappa_p(p)+y\kappa_p(u)
\tag{A.63.16}
\]
is an isomorphism: valuation and logarithm give its diagonal
matrix with entries $1,\log_p(u)$.
In these coordinates the actual Tate class is
\[
             N_p\kappa_p(p)
               +\frac{\log_p(q_{0,p})}{\log_p(u)}
                                      \kappa_p(u).
\tag{A.63.17}
\]
Thus two rational generators already span the whole local
cohomology space by the displayed map.  The simultaneous
global exponent constraints and multiplier equations for
several labels remain to be resolved; no arbitrary-support
recognition conclusion is claimed here.
